% TOP_PROBLEM: {"id": 3324, "title": "Are different notions of the crossing number the same?", "category_id": 3, "category": {"id": 3, "name": "graph_theory", "display_name": "Graph Theory", "description": "Problems involving graphs, networks, and their properties.", "slug": "graph-theory", "order_index": 3, "created_at": "2026-07-31T15:26:25.671Z"}, "status": "open", "problem_number": "OPG-37117", "set_id": 12}
% FIRST_POSTED: 2026-10-01
% ENTRY_KIND: statement_only
% UnsolvedMath ID 3324; source catalogue code OPG-37117
% !TeX program = lualatex
% Definition and sources only; this file is not an LLM solution attempt.
\documentclass[11pt,a4paper]{article}
\usepackage[margin=25mm]{geometry}
\usepackage{fontspec}
\setmainfont{lmroman10-regular.otf}[BoldFont=lmroman10-bold.otf,
 ItalicFont=lmroman10-italic.otf,BoldItalicFont=lmroman10-bolditalic.otf]
\usepackage{amsmath,amssymb,mathtools}
\usepackage{unicode-math}
\setmathfont{latinmodern-math.otf}
\AtBeginDocument{\let\setminus\smallsetminus}
\usepackage{xurl,hyperref}
\hypersetup{colorlinks=true,urlcolor=blue}
\setcounter{secnumdepth}{0}
\setlength{\parindent}{0pt}
\setlength{\parskip}{5pt}
\setlength{\emergencystretch}{3em}
\begin{document}
\section*{Are different notions of the crossing number the same?}\label{3324}
{\small UnsolvedMath ID 3324 · OPG-37117 · Graph Theory · OpenGarden}
\subsection{Definitions and mathematical statement}
Let $G$ be a finite simple graph. Consider topological plane drawings
with distinct vertex points and simple edge arcs avoiding all other
vertices. Interior intersections are finitely many transverse
crossings; no three edges cross at one point and no two edge arcs
overlap along a segment. Two edges are allowed to cross several
times. A common endpoint is not a crossing.

For a drawing $\mathcal D$, let $c(\mathcal D)$ be its number of
crossing points, and let $p(\mathcal D)$ be the number of unordered
pairs of distinct edges whose interiors cross at least once. Pairs
are counted once even if they cross many times; this includes pairs
of adjacent edges if they cross away from their common endpoint.
Define
\[
 \operatorname{cr}(G)=\min_{\mathcal D}c(\mathcal D),
 \qquad
 \operatorname{pcr}(G)=\min_{\mathcal D}p(\mathcal D).
\]
The drawings realizing the two minima may be different.

The problem asks whether $\operatorname{pcr}(G)=\operatorname{cr}(G)$
for every finite simple graph. One always has
$\operatorname{pcr}(G)\le\operatorname{cr}(G)$ because a crossing edge
pair accounts for at least one crossing point. The difficult issue
is whether allowing repeated crossings of a pair can lower the
number of different pairs that must cross.

In defining the pair parameter, one must not restrict in advance
to drawings where each edge pair crosses at most once: that would
remove precisely the phenomenon under investigation. This parameter
also differs from odd-crossing number, which counts only pairs
crossing an odd number of times.
\subsection{Short English statement}
Can minimizing the number of crossing edge pairs ever give a smaller value than minimizing the total number of crossing points?
\subsection{Sources}
\begin{itemize}
\item[S1] ULAM.AI and the UnsolvedMath Contributors, supplied problems.json, record 3324 (OPG-37117), OpenGarden collection. Catalogue identity and source transcription; CC BY 4.0.
\url{https://huggingface.co/datasets/ulamai/UnsolvedMath}.
\item[S2] Open Problem Garden, Are different notions of the crossing number the same?. Original problem and discussion; the mathematical formulation below is expanded from the supplied source record.
\url{http://www.openproblemgarden.org/op/are_different_notions_of_the_crossing_number_the_same}.
\item[S3] J. Karl and G. Tóth, A slightly better bound on the crossing number in terms of the pair-crossing number (2021), comparison of the two drawing parameters.
\url{https://arxiv.org/abs/2105.14319}.
\end{itemize}
\end{document}
